\documentclass[12pt,a4paper]{article}
\usepackage[utf8]{inputenc}
\usepackage[indonesian]{babel}
\usepackage[margin=1in]{geometry}
\usepackage{tabularx}
\usepackage{enumitem}
\usepackage{parskip}

\title{\textbf{Ekstraksi Materi: Perumusan dan Dinamika Dasar Negara Pancasila}}
\author{}
\date{}

\begin{document}
\maketitle

Sejak awal tahun 2020, bangsa Indonesia menghadapi pandemi \textit{Covid-19} yang akhirnya terjadi di seluruh Indonesia. Semua pihak ikut terlibat dalam penanganan pandemi tersebut. Keterlibatan semua pihak dalam penanganan pandemi menunjukkan bahwa bangsa Indonesia bersatu padu menghadapi semua masalah yang sedang menimpa bangsa dan negara. Hal tersebut menunjukkan bahwa masyarakat Indonesia senang bekerja sama. Sikap kerja sama dalam masyarakat menunjukkan penerapan nilai-nilai Pancasila dalam masyarakat karena Pancasila sebagai dasar negara sudah menyatu sebagai jiwa bangsa Indonesia.

Setelah pandemi \textit{Covid-19} berakhir, kalian dapat mengambil hikmah pentingnya menerapkan kebiasaan hidup sehat bagi pribadi, keluarga, maupun masyarakat. Bangsa yang berkarakter Pancasila dibangun di atas jiwa dan raga yang berkebiasaan sehat. Oleh karena itu, membutuhkan komitmen, kerja sama, dan gotong royong untuk mewujudkannya dengan menerapkan sila-sila Pancasila.

Setelah ditetapkan sebagai dasar negara, Pancasila berkaitan erat dengan dinamika pelaksanannya dalam sikap dan tindakan sehari-hari. Bagaimana rumusan dan keterkaitan sila-sila dalam Pancasila? Apa makna sila-sila Pancasila? Bagaimana kedudukan Pancasila sebagai ideologi negara dan penerapannya dalam kehidupan sehari-hari? Simaklah materi berikut untuk menjawab semua pertanyaan tersebut.

\section*{A. Rumusan Dasar Negara Indonesia Merdeka}
Sebelum membahas lebih dalam tentang Pancasila dalam sikap dan tindakan keseharian, tidak ada salahnya jika kalian mengingat kembali tentang rumusan dasar negara Indonesia merdeka. Tentunya, kalian masih ingat para tokoh anggota BPUPK dan pandangan-pandangannya, bukan? Pada sidang pertama tanggal 29 Mei 1945, Radjiman Wedyodiningrat sebagai Ketua BPUPK meminta para anggotanya memberikan pandangan tentang dasar negara Indonesia merdeka. Ketua BPUPK melontarkan pertanyaan pemantik, \textit{Negara Indonesia merdeka yang akan kita bangun itu, apa dasarnya?} Semua tokoh dalam Sidang BPUPK menyampaikan pandangannya tentang rumusan dasar negara. Ulasan tentang pendapat dan rumusan dari para tokoh telah dibahas di Kelas X. Tentunya, kalian masih ingat, bukan?

\subsection*{1. Perumusan Dasar Negara}
Pada akhir masa sidang pertama tanggal 1 Juni 1945, BPUPK sepakat untuk membentuk Panitia Kecil yang tugasnya menampung usul-usul yang masuk dan memeriksanya serta melaporkan kepada sidang pleno BPUPK. Adapun anggota panitia kecil ini terdiri atas delapan orang, yaitu sebagai berikut.
\begin{enumerate}[label=\alph*.]
    \item Ir. Sukarno (sebagai ketua)
    \item Ki Bagoes Hadikoesoemo
    \item KH Wachid Hasjim
    \item Mr. Mohammad Yamin
    \item M. Soetardjo Kartohadikoesoemo
    \item Mr. A.A. Maramis
    \item R. Otto Iskandardinata
    \item Drs. Mohammad Hatta
\end{enumerate}

Pada akhir sidang Panitia Kecil, kemudian dibentuk Panitia Kecil yang beranggotakan sembilan orang (kemudian dikenal dengan Panitia Sembilan) yang bertugas untuk menyusun rancangan Pembukaan Undang-Undang Dasar Negara Republik Indonesia yang di dalamnya termuat dasar negara. Anggota Panitia Sembilan adalah sebagai berikut.
\begin{enumerate}[label=\alph*.]
    \item Ir. Sukarno (sebagai ketua)
    \item Drs. Mohammad Hatta
    \item Ahmad Soebardjo
    \item KH Abdoel Kahar Moezakir
    \item Abikoesno Tjokrosoejoso
    \item KH Wachid Hasjim
    \item Mr. Mohammad Yamin
    \item Mr. A.A. Maramis
    \item H. Agoes Salim
\end{enumerate}

Panitia Sembilan dibentuk sebagai upaya mempertemukan pandangan antara dua golongan (kebangsaan dan Islam) mengenai dasar kenegaraan. Awalnya, terdapat kesulitan untuk mencari kecocokan paham. Namun, dapat menyepakati rumusan \textit{preambul} yang di dalamnya terdapat rumusan Pancasila. Setelah melewati proses persetujuan bersama, rumusan Pancasila 1 Juni 1945 mengalami penyempurnaan.

\vspace{0.5cm}
\textbf{Tabel 1.1 Penyempurnaan Rumusan Pancasila}
\vspace{0.2cm}

\begin{tabularx}{\textwidth}{|c|p{3.5cm}|X|X|}
\hline
\textbf{No.} & \textbf{Sila} & \textbf{Rumusan} & \textbf{Penyempurnaan} \\
\hline
1. & Ketuhanan & Dipindah dari sila terakhir ke sila pertama. & Penambahan anak kalimat \textit{dengan kewajiban menjalankan syariat Islam bagi pemeluk-pemeluknya.} \\
\hline
2. & Internasionalisme atau peri kemanusiaan & Tetap diletakkan pada sila kedua. & Penyempurnaan redaksi \textit{Kemanusiaan yang adil dan beradab.} \\
\hline
3. & Kebangsaan Indonesia & Berubah dari sila pertama menjadi sila ketiga. & Redaksinya menjadi \textit{Persatuan Indonesia.} \\
\hline
4. & Mufakat atau demokrasi & Berubah dari sila ketiga menjadi sila keempat. & Redaksinya menjadi \textit{Kerakyatan yang dipimpin oleh hikmat kebijaksanaan dalam permusyawaratan/perwakilan.} \\
\hline
5. & Kesejahteraan sosial & Berubah dari sila keempat menjadi sila kelima. & Redaksinya menjadi \textit{Keadilan sosial bagi seluruh rakyat Indonesia.} \\
\hline
\end{tabularx}
\vspace{0.5cm}

Selanjutnya, Panitia Sembilan berhasil merumuskan dan menyetujui rancangan Pembukaan UUD pada tanggal 22 Juni 1945 yang disebut dengan Piagam Jakarta. Hasil rumusan Piagam Jakarta dan berbagai usulan yang berhasil dihimpun, kemudian dilaporkan dan dibahas pada masa persidangan kedua BPUPK pada tanggal 10--17 Juli 1945. Pada sidang hari pertama (10 Juli 1945) diambil keputusan tentang bentuk negara.

Menurut Yudi Latif (2020), hasil rumusan Piagam Jakarta menimbulkan perdebatan menyangkut pencantuman tujuh kata sebagai anak kalimat dari Sila Ketuhanan dengan segala turunannya. Keberatan pencantuman tujuh kata tersebut bukan hanya datang dari golongan kebangsaan, melainkan ada variasi pandangan di kalangan golongan Islam. Bagi sebagian golongan kebangsaan, pencantuman tujuh kata yang mengandung perlakuan khusus terhadap umat Islam dirasa tidak cocok dalam suatu hukum dasar yang menyangkut warga negara secara keseluruhan.

Pada sidang tanggal 11 Juli 1945, Piagam Jakarta mendapat tanggapan dari J. Latuharhary. Dalam tanggapannya, ia merasa keberatan atas pencantuman tujuh kata, yaitu \textit{Ketuhanan dengan kewajiban menjalankan syariat Islam bagi pemeluk-pemeluknya}. Tanggapan Latuharhary menyebabkan perdebatan, namun berkat kewibawaan Sukarno maka untuk sementara waktu dapat diatasi. Dengan demikian, rumusan Piagam Jakarta (dengan tujuh kata) bertahan hingga akhir masa persidangan kedua tanggal 17 Juli 1945. Pada akhir sidangnya, BPUPK berhasil menyusun dasar negara (Pancasila) dalam Piagam Jakarta sebagai norma dasar (\textit{grundnorm}) yang menjiwai perumusan UUD sebagai aturan dasar (\textit{grund gesetze}).

Karena dianggap telah menyelesaikan tugasnya, BPUPK dibubarkan pada tanggal 7 Agustus 1945. Agenda berikutnya adalah menyiapkan dan mematangkan serta mengesahkan hal-hal penting untuk mempersiapkan kemerdekaan Indonesia.

Menurut Yudi Latif (2020), para pembentuk dasar negara Indonesia, dengan segala keragaman asal-usul agama dan keyakinan, ideologi, dan kearifan kesukuan yang melatarbelakanginya mampu melihat titik temu sebagai dasar simplisitas dan kompleksitas nilai dan keyakinan di Indonesia dalam lima prinsip Pancasila.

Sejak tanggal 18 Agustus 1945, Pancasila merupakan dasar filsafat negara dan pandangan hidup (\textit{philosofische grondslag}) dan juga pandangan dunia (\textit{weltanschauung}) bagi bangsa Indonesia. Pancasila sebagai dasar filsafat negara (\textit{philosofische grondslag}), karena mengandung unsur-unsur alasan filosofis berdirinya suatu negara dan setiap produk hukum di Indonesia harus berdasarkan nilai Pancasila. Pancasila sebagai pandangan hidup bangsa (\textit{weltanschauung}) karena mengandung unsur-unsur nilai-nilai agama, budaya, dan adat istiadat.

\subsection*{3. Nilai-Nilai Luhur dalam Perumusan dan Pengesahan Dasar Negara}
Sebagaimana kalian ketahui bahwa anggota BPUPK dan PPKI berasal dari representasi keberagaman sosial dan politik pada waktu itu. Dengan semangat kebersamaan dan persatuan, para anggota BPUPK dan PPKI melaksanakan sidang untuk menyusun dasar negara dan membentuk bangunan Negara Indonesia merdeka. Ada beberapa nilai luhur dalam perumusan dan pengesahan dasar negara dalam Sidang BPUPK dan PPKI tersebut. Nilai-nilai luhur tersebut adalah menghargai pendapat, rendah hati, kerja keras, keberanian, rela berkorban, menerima keputusan rapat, mengutamakan persatuan dan kesatuan, serta melaksanakan keputusan bersama.

\textbf{a. Menghargai Pendapat} \\
Para anggota dengan sopan dan tertib mendengarkan apa yang disampaikan oleh pimpinan sidang dan pendapat yang dikemukakan oleh anggota lainnya. Pendapat yang muncul dihargai, meskipun disampaikan secara singkat maupun panjang lebar, bahkan hingga satu jam. Penyampaian pendapat disampaikan secara sopan, jelas, tidak memotong pembicaraan, tidak memaksakan pendapat, dan mengutamakan musyawarah mufakat. Perbedaan pendapat yang ada tidak menyebabkan para anggota terpecah belah. Bila ada perbedaan pendapat maka ditempuh dengan cara damai, bukan dengan kekerasan dan paksaan.

\textbf{b. Rendah Hati} \\
Dalam perumusan dasar negara, para tokoh berdebat dan menyampaikan pendapat. Para tokoh tidak sombong dengan ego dan pandangannya masing-masing. Apabila ada pendapat yang lebih sesuai untuk kepentingan bangsa dan negara, mereka menerimanya.

\textbf{c. Kerja Keras} \\
Semua anggota BPUPK bekerja keras dengan aktif menghadiri sidang dan mengerahkan segenap pikiran dan kemampuan untuk merumuskan dasar negara. Jika terjadi perbedaan pendapat dilakukan melalui lobi-lobi secara sehat dan tidak bertentangan dengan prinsip kebenaran.

\textbf{d. Keberanian} \\
Di bawah tekanan Jepang, para tokoh berani menghadapi risiko apa pun untuk kepentingan dan masa depan bangsa.

\textbf{e. Rela Berkorban} \\
Perumusan dan pengesahan dasar negara membutuhkan pengorbanan seperti waktu, biaya, tenaga, dan sebagainya. Pengorbanan merupakan bakti kepada negara.

\textbf{f. Menerima Keputusan Rapat} \\
Persidangan atau musyawarah membahas sesuatu untuk menghasilkan kesepakatan atau keputusan. Keputusan yang diambil harus diupayakan dapat diterima dengan ikhlas dan terbuka, meskipun keputusan bersama itu tidak sesuai dengan pendapat pribadi.

\textbf{g. Mengutamakan Persatuan dan Kesatuan} \\
Meskipun berbeda pandangan, para tokoh BPUPK dan PPKI mengutamakan persatuan dan kesatuan bangsa dan negara. Contohnya dari ketidaksetujuan wakil-wakil Kristen dan Katolik atas teks dalam Piagam Jakarta. Para tokoh Islam yang berbeda pandangan dapat menerima keberatan tersebut karena lebih mementingkan persatuan dan kesatuan bangsa dan negara.

\textbf{h. Melaksanakan Keputusan Bersama} \\
Melaksanakan keputusan bersama sebagai hasil musyawarah dilakukan oleh para tokoh pendiri negara. Mereka sepakat menerima dasar negara Pancasila untuk kepentingan bangsa dan negara.

\subsection*{4. Dinamika Dasar Negara Pancasila}
Sebagaimana kalian ketahui, UUD NRI Tahun 1945 yang ditetapkan oleh PPKI pada tanggal 18 Agustus 1945 mengalami perubahan dalam perjalanan sejarah ketatanegaraan. Perubahan tersebut seperti terlihat pada tabel berikut.

\vspace{0.5cm}
\textbf{Tabel 1.4 Dinamika UUD NRI Tahun 1945}
\vspace{0.2cm}

\begin{tabularx}{\textwidth}{|c|X|X|}
\hline
\textbf{No.} & \textbf{Periode Pelaksanaan} & \textbf{UUD yang Berlaku} \\
\hline
1. & 18 Agustus 1945 -- 27 Desember 1949 & UUD NRI Tahun 1945 \\
\hline
2. & 27 Desember 1949 -- 17 Agustus 1950 & Konstitusi RIS \\
\hline
\end{tabularx}

\vspace{1cm}
\begin{center}
\fbox{
    \parbox{0.9\textwidth}{
        \textbf{Tugas 1.3} \\
        Buatlah kelompok beranggotakan tiga orang. Kemudian, diskusikan tentang beberapa nilai luhur dalam perumusan dan pengesahan dasar negara yang dilakukan oleh para pendiri negara. Beberapa nilai tersebut berhubungan dengan keputusan yang diambil untuk kepentingan bangsa dan negara. Menurut kalian, apa saja nilai-nilai tersebut yang berkaitan langsung dengan keputusan akhir tentang perumusan dan pengesahan dasar negara. Diskusikan dalam kelompok kalian dan presentasikan hasilnya di depan kelas.
    }
}
\end{center}

\end{document}
